\subsection{A complete list of global necessary conjugators}

\begin{statement}{D23.P22}{Proved}{At most 144 global conjugators before the mixed tests}
\end{statement}
For an actual multiplier-one full-GT element in the normalized
state
\[
                     \rho=0,\qquad a_2=0,
\]
and either fixed value of $a_3$, its actual permutation
$V_h$ on the degree-1344 orbit belongs to a completely
enumerated list of 144 permutations. Every permutation in
that list fixes each of the twelve orthogonal blocks,
belongs to $S^{12}$, commutes with the literal first braid
generator $A$ and the nested full twist
$Z=(ABA)^2=(AB)^3$, and has the unique four-block
restriction determined in D23.P21.

The list is exhaustive for precisely these conditions.
It consists of 18 choices on block zero and eight choices
on block one, with all other blocks determined by transport.
The two degree-72 mixed component tests of D23.C25,
including their common-word graph, leave 72 permutations
for each value of $a_3$. These remaining permutations are
necessary candidates only. In particular they are not
asserted to arise from actual GT elements or to satisfy all
remaining finite evaluations.

\begin{proof}
Write $Z=(AB)^3$. In $B_3$ this is the central full twist,
so $A$ and $Z$ commute. The standard Drinfeld action fixes
$A$ and $Z$ when $\lambda=1$. The actual covariance
D23.P18 consequently implies
\[
                         V_hA=AV_h,\qquad V_hZ=ZV_h.
 \tag{D23.P22.1}
\]
By D23.F32, the normalized actual $V_h$ belongs to the
independent simple product $S^{12}$.

The degree-12 action of $\langle A,Z\rangle$ has orbits
\[
       \{0,9\},\qquad\{1,2,3,4,5,6\},\qquad
       \{7,8,10\},\qquad\{11\}.                          \tag{D23.P22.2}
\]
On the first orbit, $A$ fixes both blocks and $Z$
interchanges them. On the second orbit, the degree-three
action of $A$ and the degree-two action of $Z$ generate a
regular $C_3\times C_2$ action. This is an exact check of
the two displayed block permutations. The other two orbits
are already fixed by D23.P21's complete local lists.

We use the following elementary transport fact. Let an
abelian group generated by commuting permutations $a,z$
act transitively on a set of blocks. A permutation $v$
fixing each block and commuting with $a,z$ is determined
by its restriction to one block. That restriction must
centralize the image of the stabilizer of that block.
Conversely every such local centralizer element extends
uniquely: transport it by any word taking the base block
to the desired block. Two choices differ by a stabilizer
element, so this is well defined. The extension commutes
with both generators by construction.

For the orbit $\{0,9\}$, the stabilizer is the image of
$\langle A,Z^2\rangle$. For the regular six-block orbit,
it is the image of $\langle A^3,Z^2\rangle$. In the
group-transversal coordinates of D23.F30 the corresponding
local centralizers in $S$ have exact orders
\[
 \begin{split}
 \left|C_S\bigl(A_0,(Z^2)_0\bigr)\right|&=18,\\
 \left|C_S\bigl((A^3)_1,(Z^2)_1\bigr)\right|&=8.
 \end{split}                                               \tag{D23.P22.3}
\]
The first local pair has element orders $(3,2)$ and the
second has orders $(4,2)$. The seed certificate constructs
the actual group $S$ as the derived subgroup of the
112-point block-stabilizer image, computes both complete
centralizers by stabilizer chains, enumerates every element,
and verifies its membership in $S$ and commutation with
both local operators.

Transport as above gives a bijection between the $18\cdot8$
seed choices and all possible extensions of the known
four-block action satisfying (D23.P22.1). The global
certificate constructs every extension as a literal
permutation on all 1344 points and checks both equations
in (D23.P22.1) on every point. It also checks the block
labels and distinguishes all 144 permutations. The
transporters belong to $\Gamma$, which normalizes the
independent factors of $S^{12}$, so every extension belongs
to $S^{12}$. This proves exhaustiveness for the stated
conditions, including the equal-sign control case $a_3=0$.

For completeness, the further mixed-component restriction
has the following actual-source interpretation. Put
$X_3=C^2$, $X_4=Z$, and $R=f(X_3,X_4)$. The same
covariance gives
\[
     V_hX_3V_h^{-1}=RX_3R^{-1},\qquad
     V_hX_4V_h^{-1}=X_4.                                  \tag{D23.P22.4}
\]
The corresponding transformation of the free pair is an
automorphism for actual gentle $f$ (it is the swapped
standard action, using the two-cycle). Hence $V_h$
normalizes the image $\langle X_3,X_4\rangle$ and permutes
its orbits. Each of the two degree-72 orbits contains 24
points in blocks $7,8,10$, on which the already determined
$V_h$ is the identity. They are therefore preserved
individually, rather than exchanged.

Each degree-72 image has order 995328, derived order
82944, and the exact marked $S_4$ quotient. At
$\rho=a_2=0$, the required marked $S_4$ value is the
identity. D23.C25 enumerates the 6912 elements of the
derived group with this marked value, and propagates
(D23.P22.4) from the known 24 points. Exactly sixteen
$R$ values and four induced permutations remain in each
component. The common-word joint image has the same order
995328 as either component, so the two values are tied by
its exact graph. In the stored four-element permutation
lists, the permitted index pairs are
\[
                       (0,0),(1,3),(2,1),(3,2).
\]
Restricting each of the 144 global permutations to both
components and imposing this graph leaves exactly 72.
Only the pairs $(0,0)$ and $(1,3)$ occur among them.
Equivalently, all eighteen block-zero seeds remain, while
four of the eight block-one seeds remain. The count is
the same for the two cubic-bit control cases. This is a
necessary restriction; other mixed and first-pair conditions
have not been omitted from the final problem by calling
these 72 permutations solutions.
\end{proof}

\paragraph{Certificates and exact scope.}
The canonical seed computation is
\path{hurwitz_seed_centralizers.py}, with JSON and NPZ
records. It independently rebuilds both centralizers and
keeps the previously recorded numbering of their elements
so that downstream indices remain stable.

The global construction and mixed restriction are
\path{hurwitz_global_conjugators.py}, with JSON and NPZ
records. The arrays \texttt{all\_V\_0}, \texttt{all\_V\_1}
contain the complete lists of 144 permutations
(each row has 1344 entries); the arrays
\texttt{mixed72\_V\_0}, \texttt{mixed72\_V\_1} contain
the 72 remaining permutations in each control case.
Their corresponding label arrays keep both seed indices
and both component-permutation indices.

The mixed inputs are the exact D23.C25 records
\path{characters/records/hurwitz_mixed72_partial.npz} and
\path{characters/records/hurwitz_mixed72_joint_probe.json},
with their code. The actual-source dependencies are
D23.P18, D23.P21, D23.F32 and the supplied two-cycle;
the group structure is D23.F30. The theorem establishes
a complete finite list that every normalized actual
conjugator must enter. It establishes neither existence
of a mixed full-GT state nor its impossibility.
